% chapters/04/section-02.tex
\section{Membaca Masukan sebagai Peristiwa}
\label{sec:membaca-masukan}

\subsection{Mengapa Cara pada Bab III Tidak Cukup}

Program jendela pada Bab~III menunggu peristiwa dengan
\code{pygame.event.wait()}. Perintah itu \textbf{menghentikan} program sampai
ada peristiwa yang terjadi. Cara tersebut memadai untuk program yang hanya
menampilkan informasi, tetapi tidak untuk permainan.

Bayangkan apa yang terjadi apabila \code{event.wait()} dipakai di dalam
siklus permainan. Ketika pemain tidak menekan apa pun dan tidak menggerakkan
tetikus, tidak ada peristiwa yang datang, sehingga program berhenti menunggu.
Akibatnya tahap pembaruan dan penayangan tidak pernah dikerjakan, dan
permainan membeku. Permainan yang benar harus tetap berjalan walaupun pemain
sedang diam, karena objek-objek di dalamnya terus bergerak dan waktu terus
berjalan.

\subsection{Antrean Peristiwa dan Cara Mengosongkannya}

Pygame mengumpulkan seluruh peristiwa yang terjadi ke dalam sebuah
\textbf{antrean} (\istilah{queue}). Peristiwa masuk ke antrean ketika terjadi,
dan menunggu di sana sampai ada yang membacanya. Perintah
\code{pygame.event.get()} mengambil \textbf{seluruh} isi antrean sekaligus dan
mengosongkannya, tanpa menunggu peristiwa baru datang.

Karena itulah \code{event.get()} dipakai di dalam siklus, sedangkan
\code{event.wait()} tidak. \code{event.get()} mengembalikan
daftar, sehingga hasilnya dibaca dengan perulangan.

\begin{kodekeluaran}
for event in pygame.event.get():
    if event.type == pygame.QUIT:
        running = False
    elif event.type == pygame.KEYDOWN:
        if event.key == pygame.K_ESCAPE:
            running = False
\end{kodekeluaran}

\begin{peringatanbox}
Antrean peristiwa harus dikosongkan setiap bingkai. Apabila
\code{pygame.event.get()} tidak pernah dipanggil, antrean terus bertambah
panjang dan sistem operasi menganggap program tidak menanggapi masukan. Pada
Windows, jendela program semacam itu akan ditandai sebagai ``tidak
menanggapi'' dan dapat ditutup paksa oleh sistem.
\end{peringatanbox}

\subsection{Jenis Peristiwa yang Lazim}

Setiap peristiwa memiliki jenis (\code{event.type}) dan, untuk sebagian
peristiwa, keterangan tambahan. Peristiwa \code{KEYDOWN} misalnya memuat
\code{event.key}, yaitu tombol yang ditekan. Jenis peristiwa yang paling
sering dipakai pada Blok~I dan Blok~II diperlihatkan pada
Tabel~\ref{tab:jenis-peristiwa}.

\begin{table}[htbp]
\centering
\caption{Jenis peristiwa yang paling sering dipakai pada Blok I dan Blok II}
\label{tab:jenis-peristiwa}
{\setlength{\tabcolsep}{4pt}\footnotesize%
\begin{tabularx}{\textwidth}{@{}L{4.0cm} L{4.2cm} Y@{}}
\toprule
\headrow \textbf{Jenis peristiwa} & \textbf{Keterangan tambahan} & \textbf{Dipakai untuk} \\
\midrule
\code{QUIT} & tidak ada &
Menutup jendela: pemain meminta program berakhir. \\
\addlinespace
\code{KEYDOWN} & \code{event.key}, \code{event.mod} &
Tombol ditekan. Peristiwa ini terjadi sekali saja untuk setiap penekanan, tidak
berulang selama tombol ditahan. \\
\addlinespace
\code{KEYUP} & \code{event.key} &
Tombol dilepas. Berguna untuk tindakan yang berakhir ketika tombol
dilepaskan. \\
\addlinespace
\code{MOUSEBUTTONDOWN} & \code{event.pos}, \code{event.button} &
Tombol tetikus ditekan, beserta letak penunjuk saat itu. \\
\addlinespace
\code{MOUSEMOTION} & \code{event.pos}, \code{event.rel} &
Penunjuk tetikus bergeser. \\
\bottomrule
\end{tabularx}}
\tabsource{Sumber: dokumentasi Pygame \mbox{[R3]}.}
\end{table}

\subsection{Peristiwa Berbeda dari Keadaan Tombol}

Satu hal yang sering membingungkan pemula perlu ditegaskan di sini.
\code{event.get()} melaporkan \textbf{peristiwa}, yaitu sesuatu yang baru saja
terjadi. Peristiwa \code{KEYDOWN} terjadi sekali pada saat tombol mulai
ditekan, dan tidak terjadi lagi selama tombol itu tetap ditahan.

Untuk gerak yang berlangsung terus selama tombol ditahan---misalnya karakter
yang berjalan ke kanan---diperlukan cara lain, yaitu memeriksa
\textbf{keadaan} seluruh tombol pada saat itu dengan
\code{pygame.key.get_pressed()}. Perbedaan keduanya akan Anda pakai bersama
pada Bab~V, dan memilih cara yang keliru merupakan salah satu penyebab
gerakan yang tersendat: karakter yang hanya bergerak satu piksel setiap kali
tombol ditekan \cite{pygame-docs}.

Sebagai pedoman sederhana, gunakan peristiwa untuk tindakan yang terjadi
sekali (melompat, menembak, membuka menu, mengakhiri program), dan gunakan
keadaan tombol untuk tindakan yang berlangsung selama tombol ditahan
(berjalan, mengarahkan, mempercepat).

\begin{catatabox}
Selain peristiwa bawaan Pygame, Anda dapat membuat peristiwa sendiri dan
menaruh peristiwa itu ke antrean dengan \code{pygame.event.post()}. Cara itu
dipakai pada program buku ini pada pilihan \code{--auto-close}, yaitu untuk
mengirim peristiwa \code{QUIT} secara sengaja agar program dapat diuji tanpa
jendela. Cara yang sama berguna nanti untuk menyampaikan kejadian di dalam
permainan, misalnya ``skor bertambah'', kepada bagian program yang lain.
\end{catatabox}
